% =====================================================================================
% AGRADECIMIENTOS, RESUMEN (Revisión 4) Y SUMMARY
% Escribe tu texto FUERA de los recuadros \begin{guia} ... \end{guia}.
% =====================================================================================

\clearpage
\ClearWallPaper
\pagestyle{fancy}

\section*{\centering Agradecimientos}
\addcontentsline{toc}{section}{Agradecimientos}
% >>> ESCRIBE AQUÍ tus agradecimientos (opcional, breve) <<<


\clearpage
\section*{\centering Resumen} % audit:req=resumen
\addcontentsline{toc}{section}{Resumen}

\begin{guia}[Guía: Resumen (Revisión 4)]
\textbf{¿Para qué sirve?} Es la síntesis de todo el trabajo para quien no leerá el documento completo. Se escribe al
final, cuando ya tienes los resultados.

\textbf{Debe responder, en un solo párrafo de 150 a 250 palabras:} ¿cuál es el contexto y el problema? ¿cuál es el
objetivo? ¿qué se hizo (método o solución)? ¿cuáles fueron los resultados principales (con datos)? ¿cuál es la
conclusión o aportación?

\textbf{Reglas:} sin citas, sin figuras, sin siglas sin definir. Cierra con \textbf{4 a 6 palabras clave} (en
minúsculas, separadas por comas).

\textbf{Extensión mínima:} 120 palabras.
\end{guia}

% >>> ESCRIBE AQUÍ el resumen (un párrafo) <<<

\textbf{Palabras clave:} palabra1, palabra2, palabra3, palabra4.


\clearpage
\section*{\centering Summary}
\addcontentsline{toc}{section}{Summary}

\begin{guia}[Guía: Summary (Revisión 4)]
Traducción al inglés del Resumen (mismo contenido y estructura), con \textbf{Keywords}. Revisa la gramática; evita
traducción automática literal.
\end{guia}

% >>> WRITE HERE the English summary <<<

\textbf{Keywords:} keyword1, keyword2, keyword3, keyword4.
